\section{Data Mixtures}
\label{app:data-mixtures}

\textbf{Language lists.} We generate the language lists from the FineWeb-2 per-language byte counts and from the benchmarks available in the evaluation harness. \emph{Scheme A} ranks the FineWeb-2 subsets by training-split UTF-8 bytes, excluding the unidentified-language subsets, and takes the top $K-1$. The lists are nested. The $K=2$ list (Russian) is a subset of the $K=8$ list (Russian, Chinese, German, Japanese, Spanish, French, Italian), which is a subset of the $K=15$, $K=30$ and $K=50$ lists. We also generated a 100-language list with the same rule, but we did not train it, as explained below. In that list, the eight subsets among the top 99 that have no benchmark in the harness were replaced by the next subsets by bytes with at least two benchmark families. \emph{Scheme B} builds diversity-first lists at $K \in \{8, 15, 30\}$ within the top 49. It first takes the largest subset of each script not yet covered, then of each language family not yet covered, then the remaining subsets by bytes. At $K=8$, the scheme-B list spans seven scripts against four for scheme A. At $K=30$, it spans 14 scripts and 12 families against 10 and 10. At $K=2$ there is no list to diversify. There, the scheme-B and scheme-C builds (\emph{ZH} and \emph{ES}) replace Russian with Chinese or Spanish. At $K=1$ there is no second language. There, the scheme-B and scheme-C builds replace the English corpus instead. Scheme B (\emph{DCLMP}) uses the same DCLM pipeline without its educational-quality filter. Scheme C (\emph{FWEB}) uses FineWeb restricted to the crawls up to 2022, which is DCLM's own Common Crawl window.
Table~\ref{tab:app-languages} lists the 99 FineWeb-2 languages of the sweep. For each language, it gives the language setting at which the language enters the scheme-A and scheme-B lists, and the number of benchmark families that the harness offers for it.

\begin{table*}[p]
\centering
\tiny
\setlength{\tabcolsep}{2.5pt}
\begin{minipage}[t]{0.49\textwidth}\centering
\begin{tabular}{llllrrr}
\toprule
Subset & Language & Script & Family & $K_A$ & $K_B$ & Fam. \\
\midrule
dclm & English & Latn & Indo-European & 1 & 1 & 25 \\
rus\_Cyrl & Russian & Cyrl & Indo-European & 2 & 8 & 14 \\
cmn\_Hani & Mandarin Chinese & Hani & Sino-Tibetan & 8 & 2 & 16 \\
deu\_Latn & German & Latn & Indo-European & 8 & 8 & 13 \\
jpn\_Jpan & Japanese & Jpan & Japonic & 8 & 8 & 10 \\
spa\_Latn & Spanish & Latn & Indo-European & 8 & 30 & 18 \\
fra\_Latn & French & Latn & Indo-European & 8 & 30 & 13 \\
ita\_Latn & Italian & Latn & Indo-European & 8 & 30 & 12 \\
por\_Latn & Portuguese & Latn & Indo-European & 15 & 30 & 11 \\
pol\_Latn & Polish & Latn & Indo-European & 15 & 30 & 8 \\
nld\_Latn & Dutch & Latn & Indo-European & 15 & 30 & 8 \\
ind\_Latn & Indonesian & Latn & Austronesian & 15 & 30 & 10 \\
vie\_Latn & Vietnamese & Latn & Austro-Asiatic & 15 & 30 & 13 \\
fas\_Arab & Persian & Arab & Indo-European & 15 & 8 & 7 \\
arb\_Arab & Standard Arabic & Arab & Afro-Asiatic & 15 & 30 & 14 \\
tur\_Latn & Turkish & Latn & Turkic & 30 & 30 & 9 \\
tha\_Thai & Thai & Thai & Kra-Dai & 30 & 8 & 4 \\
ukr\_Cyrl & Ukrainian & Cyrl & Indo-European & 30 & 30 & 9 \\
ell\_Grek & Modern Greek (1453 onward) & Grek & Indo-European & 30 & 8 & 9 \\
kor\_Hang & Korean & Hang & Koreanic & 30 & 15 & 7 \\
ces\_Latn & Czech & Latn & Indo-European & 30 & 30 & 6 \\
swe\_Latn & Swedish & Latn & Indo-European & 30 & 30 & 9 \\
hun\_Latn & Hungarian & Latn & Uralic & 30 & 30 & 9 \\
ron\_Latn & Romanian & Latn & Indo-European & 30 & 30 & 6 \\
nob\_Latn & Norwegian Bokmål & Latn & Indo-European & 30 & -- & 3 \\
dan\_Latn & Danish & Latn & Indo-European & 30 & -- & 7 \\
bul\_Cyrl & Bulgarian & Cyrl & Indo-European & 30 & -- & 7 \\
fin\_Latn & Finnish & Latn & Uralic & 30 & -- & 7 \\
hin\_Deva & Hindi & Deva & Indo-European & 30 & 15 & 13 \\
ben\_Beng & Bengali & Beng & Indo-European & 30 & 15 & 9 \\
slk\_Latn & Slovak & Latn & Indo-European & 50 & -- & 7 \\
heb\_Hebr & Hebrew & Hebr & Afro-Asiatic & 50 & 15 & 7 \\
lit\_Latn & Lithuanian & Latn & Indo-European & 50 & -- & 7 \\
bos\_Latn & Bosnian & Latn & Indo-European & 50 & -- & 1 \\
slv\_Latn & Slovenian & Latn & Indo-European & 50 & -- & 3 \\
ekk\_Latn & Standard Estonian & Latn & Uralic & 50 & -- & 7 \\
cat\_Latn & Catalan & Latn & Indo-European & 50 & -- & 8 \\
tam\_Taml & Tamil & Taml & Dravidian & 50 & 15 & 9 \\
hrv\_Latn & Croatian & Latn & Indo-European & 50 & -- & 7 \\
lvs\_Latn & Standard Latvian & Latn & Indo-European & 50 & -- & 1 \\
zsm\_Latn & Standard Malay & Latn & Austronesian & 50 & -- & 6 \\
azj\_Latn & North Azerbaijani & Latn & Turkic & 50 & -- & 5 \\
srp\_Cyrl & Serbian & Cyrl & Indo-European & 50 & -- & 8 \\
kat\_Geor & Georgian & Geor & Kartvelian & 50 & 15 & 6 \\
npi\_Deva & Nepali (individual language) & Deva & Indo-European & 50 & -- & 8 \\
mar\_Deva & Marathi & Deva & Indo-European & 50 & -- & 7 \\
mal\_Mlym & Malayalam & Mlym & Dravidian & 50 & 15 & 7 \\
kaz\_Cyrl & Kazakh & Cyrl & Turkic & 50 & -- & 6 \\
urd\_Arab & Urdu & Arab & Indo-European & 50 & -- & 7 \\
als\_Latn & Tosk Albanian & Latn & Indo-European & 50 & -- & 5 \\
\bottomrule
\end{tabular}
\end{minipage}\hfill
\begin{minipage}[t]{0.49\textwidth}\centering
\begin{tabular}{llllrrr}
\toprule
Subset & Language & Script & Family & $K_A$ & $K_B$ & Fam. \\
\midrule
mkd\_Cyrl & Macedonian & Cyrl & Indo-European & val & -- & 4 \\
tel\_Telu & Telugu & Telu & Dravidian & val & -- & 7 \\
kan\_Knda & Kannada & Knda & Dravidian & val & -- & 4 \\
mya\_Mymr & Burmese & Mymr & Sino-Tibetan & val & -- & 2 \\
guj\_Gujr & Gujarati & Gujr & Indo-European & val & -- & 5 \\
bel\_Cyrl & Belarusian & Cyrl & Indo-European & val & -- & 3 \\
isl\_Latn & Icelandic & Latn & Indo-European & val & -- & 3 \\
khm\_Khmr & Khmer & Khmr & Austro-Asiatic & val & -- & 1 \\
khk\_Cyrl & Halh Mongolian & Cyrl & Mongolic & val & -- & 1 \\
fil\_Latn & Filipino & Latn & Austronesian & val & -- & 3 \\
ary\_Arab & Moroccan Arabic & Arab & Afro-Asiatic & val & -- & 14 \\
afr\_Latn & Afrikaans & Latn & Indo-European & val & -- & 1 \\
hye\_Armn & Armenian & Armn & Indo-European & val & -- & 6 \\
sin\_Sinh & Sinhala & Sinh & Indo-European & val & -- & 3 \\
glg\_Latn & Galician & Latn & Indo-European & val & -- & 6 \\
uzn\_Cyrl & Northern Uzbek & Cyrl & Turkic & val & -- & 3 \\
pan\_Guru & Panjabi & Guru & Indo-European & val & -- & 2 \\
ory\_Orya & Odia & Orya & Indo-European & val & -- & 1 \\
uzn\_Latn & Northern Uzbek & Latn & Turkic & val & -- & 3 \\
kir\_Cyrl & Kirghiz & Cyrl & Turkic & val & -- & 4 \\
eus\_Latn & Basque & Latn & Language isolate & val & -- & 10 \\
lat\_Latn & Latin & Latn & Indo-European & val & -- & 1 \\
tgk\_Cyrl & Tajik & Cyrl & Indo-European & val & -- & 1 \\
swh\_Latn & Swahili (individual language) & Latn & Niger-Congo & val & -- & 6 \\
arz\_Arab & Egyptian Arabic & Arab & Afro-Asiatic & val & -- & 14 \\
nno\_Latn & Norwegian Nynorsk & Latn & Indo-European & val & -- & 1 \\
cym\_Latn & Welsh & Latn & Indo-European & val & -- & 1 \\
amh\_Ethi & Amharic & Ethi & Afro-Asiatic & val & -- & 4 \\
pbt\_Arab & Southern Pashto & Arab & Indo-European & val & -- & 1 \\
ckb\_Arab & Central Kurdish & Arab & Indo-European & val & -- & 2 \\
ars\_Arab & Najdi Arabic & Arab & Afro-Asiatic & val & -- & 14 \\
gle\_Latn & Irish & Latn & Indo-European & val & -- & 1 \\
lao\_Laoo & Lao & Laoo & Kra-Dai & val & -- & 1 \\
srp\_Latn & Serbian & Latn & Indo-European & val & -- & 8 \\
som\_Latn & Somali & Latn & Afro-Asiatic & val & -- & 2 \\
bod\_Tibt & Tibetan & Tibt & Sino-Tibetan & val & -- & 1 \\
mlt\_Latn & Maltese & Latn & Afro-Asiatic & val & -- & 1 \\
asm\_Beng & Assamese & Beng & Indo-European & val & -- & 2 \\
snd\_Arab & Sindhi & Arab & Indo-European & val & -- & 2 \\
kmr\_Latn & Northern Kurdish & Latn & Indo-European & val & -- & 1 \\
uig\_Arab & Uighur & Arab & Turkic & val & -- & 2 \\
plt\_Latn & Plateau Malagasy & Latn & Austronesian & val & -- & 2 \\
kin\_Latn & Kinyarwanda & Latn & Niger-Congo & val & -- & 2 \\
jav\_Latn & Javanese & Latn & Austronesian & val & -- & 2 \\
xho\_Latn & Xhosa & Latn & Niger-Congo & val & -- & 1 \\
hat\_Latn & Haitian & Latn & Creole & val & -- & 2 \\
fao\_Latn & Faroese & Latn & Indo-European & val & -- & 2 \\
zul\_Latn & Zulu & Latn & Niger-Congo & val & -- & 2 \\
ibo\_Latn & Igbo & Latn & Niger-Congo & val & -- & 2 \\
sot\_Latn & Southern Sotho & Latn & Niger-Congo & val & -- & 1 \\
\bottomrule
\end{tabular}
\end{minipage}
\caption{The languages of the sweep. $K_A$ gives the smallest language setting whose scheme-A (resource-ranked) list contains the subset. $K_B$ gives the smallest trained setting of a scheme-B build that contains it (the Chinese swap at $K=2$, the diversity-first lists at $K \in \{8, 15, 30\}$). The scheme-C build at $K=2$ trains Spanish. The lists are nested, so a language is also trained at every larger setting of its scheme. The 50 languages marked val are in no training mixture. They enter only the shared validation set. Fam.\ counts the benchmark families of the pretraining suite (reformulated variants excluded) that the harness offers for the language. English is the DCLM half of every mixture. Family names are shortened to their top-level group.}
\label{tab:app-languages}
\end{table*}



\textbf{Mixture composition.} Every multilingual setting blends English and FineWeb-2 50/50 at training time, with the blend weights of the Megatron data loader. The blend uses one English dataset and one FineWeb-2 dataset per (scheme, setting). Within the FineWeb-2 half, language $i$ receives a share $q_i \propto p_i^{1/T}$, where $p_i$ is its estimated token proportion in the corpus. Scheme A allocates at $T=1$, in proportion to the corpus. The FineWeb-2 byte distribution is extreme. On the planned 100-language list at $T=1$, the largest language would take 28\% of the multilingual half and the smallest 0.0017\%. At the 90M rung, 66 of the 99 languages would receive under 10M tokens. The flattened scheme AT3 allocates at $T=3$, which would lift the median language of that build from 90M to 373M tokens. We dropped the $K=100$ setting from the grid. Instead, we train the temperature intervention on the scheme-A lists at $K=50$ (both architectures). Where flattening starves no language, we also train it at $K=30$ and $K=15$ (deep only). At these two settings, the smallest language keeps 343M and 1.2B tokens at $T=1$ under the 1.7B budget.

\textbf{Build sizing and repetition.} Each FineWeb-2 dataset is built to half of the largest budget that trains on it, plus 10\% headroom. This gives 92B tokens for the 1.7B rung and 165B at the settings where the 3B rung trains. The ZH build (52B) was sized when that scheme stopped at 1B, and its 1.7B cell reads it unchanged. The English dataset is built once, to 184B tokens. The builder never repeats a document, so a language that runs out of data gives a smaller build. Smaller models draw a shuffled fraction of the same build, so within a setting all sizes see the same language proportions. We record two departures from a single pass and carry them into the analysis. First, no single FineWeb-2 subset can feed the 1.7B rung at $K=2$. The Russian build has 72.8B tokens, against the 83.6B that a 1.7B model draws from the multilingual half, so the 1.7B model at $K=2$ sees Russian about 1.15 times. The Chinese build (52.0B tokens) and the Spanish build (23.9B, all the Spanish that the source holds) are smaller still. At 1.7B, Chinese is seen 1.61 times and Spanish 3.51 times (0.65, 0.91 and 1.98 times for Russian, Chinese and Spanish at 1B). Repetition stays below 1.7 epochs across the Russian and Chinese ladders. For Spanish, it rises from 0.7 at 350M to 3.5 at 1.7B, so any comparison of the $K=2$ schemes states these epoch counts. Second, the $K=15$ and $K=50$ builds of scheme A and the $K=15$ build of scheme B were first built to 52B. They were extended to 92B when the 1.7B rung was added at those settings. The extension continues the token stream of each language from the same starting point. The six 1.7B cells there therefore read a superset of what their proxies read, extended with later CommonCrawl dumps.\footnote{Verified byte for byte: every language section of each 52B build is a prefix of the corresponding 92B build.}

\textbf{Validation set.} One fixed validation set is shared by every model and every scheme. It covers the 99 FineWeb-2 languages of the $K=100$ list plus English. For each language, it takes the first 5M tokens of the first parquet file of that language. It takes at most 30\% of the rows of the file, so that single-file languages keep training data. The build records the token count of each language, its UTF-8 byte count and the number of leading rows assigned to validation. Every training build reads this manifest and skips exactly those rows, so training and validation are disjoint by construction. Since there is one validation build for the whole sweep, bits-per-byte is comparable across sizes, language settings and schemes.

